%%%%%%%%%%%%%%%%%%%%%%% CHAPTER - 1 %%%%%%%%%%%%%%%%%%%%
\chapter{Introduction} 


%%%%%%%%%%%%%%%%%%%%%%%%%%%%
%%%%%%%%%%%%%%%%%%%%%%%%%%%%

% 
%%%%%%%%%%%%%%%%%%%%%%%%%%%%%%%%%%%%%%%%%%%%%%%%%%%%%%%%%%%%%%%%
%%%%%%%%%%%%%%%%%%%%%%%%%%%%%%%%%%%%%%%%%%%%%%%%%%%%%%%%%%%%%%%%

\section{Strong Interaction}

The strong interaction is one of the four fundamental forces of nature and is responsible for binding quarks into hadrons and, at a larger scale, binding protons and neutrons together to form atomic nuclei. It is described within the Standard Model of particle physics by Quantum Chromodynamics (QCD), a non-Abelian gauge theory built on the color charge carried by quarks and gluons. Unlike the electromagnetic interaction, QCD exhibits two defining features: asymptotic freedom, whereby quarks behave as nearly free particles at short distances (high energies), and confinement, whereby the interaction grows so strong at larger distances that free quarks and gluons are never observed in isolation. This interplay makes the strong interaction uniquely challenging to study, since perturbative techniques that work well at high energies break down in the low-energy regime relevant to nuclear physics.

At the energy scales relevant to nuclear structure, the strong interaction is more conveniently described not by quarks and gluons directly, but by an effective interaction between nucleons (protons and neutrons), historically modeled through meson-exchange theories and, more recently, through chiral effective field theory. This effective nucleon-nucleon interaction retains the essential residual effects of the underlying QCD dynamics and provides the foundation upon which nuclear structure and, ultimately, the physics of dense stellar matter can be understood. Establishing a reliable and consistent description of the strong interaction across these scales is therefore a prerequisite for the topics discussed in the remainder of this thesis.

\section{Nucleus}

The atomic nucleus is a self-bound, many-body quantum system governed by the residual strong interaction between its constituent nucleons. Despite decades of theoretical and experimental progress, a complete and unified description of nuclear structure remains an open problem, owing to the complexity of the nuclear force, the many-body nature of the system, and the wide range of phenomena nuclei exhibit — from simple shell-like behavior in light and near-magic nuclei to collective excitations, deformation, and clustering in heavier or exotic systems.

Nuclear models such as the shell model, mean-field approaches, and ab initio many-body methods each capture different aspects of nuclear behavior, and their predictions become increasingly important when extrapolated toward regions of the nuclear chart that are difficult or impossible to probe experimentally, such as extremely neutron-rich matter. Understanding how nucleons organize themselves under the strong force, how nuclear matter responds to changes in density and isospin asymmetry, and how the nuclear equation of state (EOS) behaves under extreme conditions are central questions that connect terrestrial nuclear physics to astrophysical phenomena. This connection forms the conceptual bridge to the final subject of this thesis.

\section{Neutron Star}

Neutron stars represent one of the most extreme environments in which the strong interaction and nuclear matter can be studied, offering a natural laboratory that is entirely inaccessible on Earth. Formed in the aftermath of core-collapse supernovae, neutron stars compress roughly a solar mass of matter into a radius of only about ten kilometers, reaching central densities several times that of ordinary atomic nuclei. Under these conditions, matter is expected to transition through a sequence of exotic phases — from a neutron-rich outer crust, through increasingly dense nuclear matter, to a core whose composition remains uncertain and may include hyperons, meson condensates, or deconfined quark matter.

The properties of neutron stars — their mass-radius relationship, maximum mass, cooling behavior, and tidal deformability during binary mergers — are directly governed by the nuclear equation of state, which is in turn rooted in the strong interaction between nucleons. Recent multi-messenger observations, including precise pulsar mass measurements and gravitational-wave signals from neutron star mergers such as GW170817, have opened an unprecedented observational window into this regime, allowing theoretical predictions derived from nuclear physics to be directly confronted with astrophysical data. This convergence of nuclear theory and astrophysical observation motivates the work presented in this thesis: to use our understanding of the strong interaction and nuclear structure to constrain the properties and internal composition of neutron stars.



\section{Conventions}
\begin{color}{blue}
	\begin{itemize}
		\item Greek indices run from $0$ to $3$ and denote space--time components (with
		$0$ being the time component), Latin indices run from $1$ to $3$ and denote
		spatial components. We use the Einstein sum convention for these indices.
		
		\item Latin indices in roman typeset, e.g.\ $\mathrm{n,p,x,\ldots}$, are constituent
		indices. We place them liberally upstairs or downstairs in order to avoid cluttering
		and they do not obey any kind of summation convention.
		
		\item The signature of the metric is $(-,+,+,+)$ which means that time-like vectors
		have a negative length.
		
		\item Unless stated otherwise, we carry out our calculations in units in which
		$G=c=M_\odot=1$. Whenever appropriate, we quote physical quantities in
		cgs-units; occasionally we use MeV.
		

		

		

	\end{itemize}
	
\end{color}

We plan the thesis in the following fashion. In Chapter \ref{C2}, we introduce the theoretical framework for neutron star 





 the quantum mechanical approach for constructing the formalism of neutrino oscillations. We provide fundamental insights into two-flavor and three-flavor neutrino oscillations, both in vacuum and Earth's matter. Furthermore, we conduct an in-depth exploration of the degeneracies associated with the neutrino oscillation parameters and highlight the prevailing tensions among various neutrino oscillation experiments. Chapter \ref{C3} presents a comprehensive overview of the evolution of long-baseline experiments, tracing their trajectory from inception to their current advancements. In Chapter \ref{C4}, we illustrate how the contemporary and next-generation long-baseline (LBL) neutrino experiments hold immense potential in probing the deviation of the atmospheric mixing angle ($\theta_{23}$) from maximal mixing, resolving the octant ambiguity ($\theta_{23}>45^\circ$ or $<45^\circ$), and significantly enhancing the precision of the oscillation parameters $\theta_{23}$ and $\Delta m^2_{31}$. We particularly emphasize the role of the upcoming Deep Underground Neutrino Experiment (DUNE) in surpassing present measurements within the three-neutrino framework. Chapter \ref{C5} delves into the synergistic interplay between DUNE and T2HK (Tokai-to-Hyper-Kamiokande) in resolving degeneracies among neutrino oscillation parameters. We investigate how the combined sensitivity of these experiments strengthens key aspects of neutrino oscillation phenomenology, including the establishment of non-maximal $\theta_{23}$, exclusion of the incorrect octant, and precision measurements of the 2-3 oscillation parameters of the PMNS matrix. These pursuits are particularly compelling as neutrino oscillation research transitions into the precision era. Additionally, we demonstrate how the complementary features of DUNE and T2HK facilitate an enhanced physics reach, achieving substantial sensitivity at lower exposures compared to their individual baseline configurations. Lastly, in Chapter \ref{C6}, we explore how the next-generation LBL experiments can uncover signatures of the flavor-dependent long-range neutrino interactions within the neutrino sector. Specifically, we analyze how these experiments can probe the presence of the gauged $(L_e - L_\mu)$, $(L_e - L_\tau)$, and $(L_\mu - L_\tau)$ symmetries through the aforementioned oscillation phenomenology, thereby establishing compelling evidence for such interactions with a high degree of confidence. Finally, Chapter \ref{C7} ends this thesis with some insightful concluding remarks, encapsulating the key takeaways, and their importance in the neutrino oscillation frontier.


%\vspace{20pt}
%%%%%%%%%%%%%%%%%%%%%%%%%%%%
%\noindent \textbf{Arrangement of the Dissertation:}
%%%%%%%%%%%%%%%%%%%%%%%%%%%%